\section{Ejecución e instalación de componentes}
\begin{itemize}
    \item[1)] \textbf{Ejecutar MadGraph}\\
    Dentro de la carpeta de MadGraph (MG5\_aMC\_v2\_9\_27) abrir la terminal e ingresar al entorno conda MG5
\begin{lstlisting}
> conda activate MG5
(MG5)> ./bin/mg5_aMC
MG5_aMC>
\end{lstlisting}
    donde \textit{MG5\_aMC>} es el entorno de consola de MadGraph.
    \item[2)] \textbf{Instalación de Pythia 8 (shower)}\\
    En la consola de MadGraph
\begin{lstlisting}
MG5_aMC> install pythia8
\end{lstlisting}
    Durante la instalación se preguntará si se desa instalar \textit{LHAPDF6}, siendo la mejor opción el rechazar la instalación de \textit{LHAPDF6}, debido a errores que puede presentar.
    \item[3)] \textbf{Instalación de interfaz MadGraph-Pythia}\\
    Después de la instalación de Pythia 8 ejecutra la siguiente linea de comando en el entorno de MadGraph
\begin{lstlisting}
MG5_aMC> install mg5amc_py8_interface
\end{lstlisting}
    \item[4)] \textbf{Instalación de Delphes (detector)}\\
    Dentro del entorno de MadGraph
\begin{lstlisting}
MG5_aMC> install Delphes
\end{lstlisting}
    la instalación puede requerir de compiladores, \textit{make} y \textit{cmake}, los cuales ya estan incluidos en el entorno MG5.
    \item[5)] \textit{Activar la conversión de modelos}\\
    Dentro del entorno de MadGraph
\begin{lstlisting}
MG5_aMC> set auto_convert_model T
MG5_aMC> exit
\end{lstlisting}
    \item[6)] \textbf{Configuración de \textit{PYTHIA8DATA}}\\
    Fuera de MadGraph, se define la ruta al directorio \textit{xmldoc},
\begin{lstlisting}
# En caso de que la carpeta de MadGraph se encuentre en la ruta:
# /home/usuario/Descargas/madgraph/MG5\_aMC\_v2\_9\_27

# Ejecutando en la terminal
(MG5)> export PYTHIA8DATA=/home/usuario/Descargas/MG5_aMC_v2_9_27/HEPTools/pythia8/share/Pythia8/xmldoc

# Comprobando la ubicacion
(MG5)> echo "$PYTHIA8DATA"
\end{lstlisting}
    \item[7)] \textbf{Añadir al archivo ~/.bashrc}\\
    Para no ejecutar el export en cada terminal, se agrega la dirección al archivo \textit{~/.bashrc} y posteriormente se recarga la configuración
\begin{lstlisting}
# Importando la direccion de los archivos xmldoc
(MG5)> echo 'export PYTHIA8DATA=/home/usuario/Descargas/madgraph/MG5_aMC_v2_9_27/HEPTools/pythia8/share/Pythia8/xmldoc' >> ~/.bashrc

# Recargando la configuracion
(MG5)> source ~/.bashrc
\end{lstlisting}
    \item[8)] \textbf{Instalación de compilador Fortran (gfortran)}\\
    Dentro del entorno conda MG5 
\begin{lstlisting}
(MG5)> conda install -c conda-forge gfortran_linux-64

# Verificando la version
(MG5)> gfortran --version

# Buscando el directorio donde se instalo el compilador
(MG5)> find "$CONDA_PREFIX/bin" -maxdepth 1 -name "*gfortran*"
<out> /home/usuario/anaconda3/envs/MG5/bin/x86_64-conda-linux-gnu-gfortran
\end{lstlisting}
    \item[9)] \textbf{Configuración automática de Fortran en el entorno MG5}
    \begin{lstlisting}
# Creacion del directorio de configuracion
(MG5)> mkdir -p "$CONDA_PREFIX/etc/conda/activate.d"

# Creando archivo nano para la configuracion de gfortran
(MG5)> nano "$CONDA_PREFIX/etc/conda/activate.d/mg5_compiler.sh"

# Dentro del archivo nano 
(mg5_compiler.sh)> export FC="$CONDA_PREFIX/bin/x86_64-conda-linux-gnu-gfortran"  # ubicacion del compilador gfortran
(mg5_compiler.sh)> export F77="$FC"
(mg5_compiler.sh)> export F90="$FC"

# Guardando y saliendo del archivo nano
(mg5_compiler.sh)> Ctrl + O
(mg5_compiler.sh)> Enter
(mg5_compiler.sh)> Ctrl + X
\end{lstlisting}
    Verificando la configuración
    \begin{lstlisting}
# Desactivando y volviendo a activar el entorno conda MG5, para actualizar los cambios
(MG5)> conda deactivate
     > conda activate MG5

# Comprobando que esta activo el compilador de fortran
(MG5)> echo "$FC"
(MG5)> "$FC" --version
\end{lstlisting}
\end{itemize}